\documentclass[11pt,a4paper]{article}

\usepackage[T1]{fontenc}
\usepackage[utf8]{inputenc}
\usepackage{lmodern}
\usepackage{amssymb,amsmath,amsthm,mathtools}
\usepackage[margin=2.5cm]{geometry}
\usepackage{parskip}
\usepackage[hidelinks]{hyperref}
\usepackage{microtype}

% Theorem environments (arXiv / journal style)
\theoremstyle{plain}
\newtheorem{theorem}{Theorem}[section]
\newtheorem{proposition}[theorem]{Proposition}
\newtheorem{lemma}[theorem]{Lemma}
\newtheorem{corollary}[theorem]{Corollary}

\theoremstyle{definition}
\newtheorem{definition}[theorem]{Definition}
\newtheorem{example}[theorem]{Example}
\newtheorem{remark}[theorem]{Remark}

% Notation
\newcommand{\Type}{\mathbf{Type}}
\newcommand{\Ok}{\mathrm{Ok}}
\newcommand{\localOk}{\mathrm{localOk}}

\title{Trust Boundaries on Registry Trees:\\
  Duplicate-Stability of the Boundary Rule}
\author{Extracted from a machine-checked development}
\date{}

\begin{document}
\maketitle

\begin{abstract}
We treat trust-labelled, finitely branching rose trees as a directed
container and formalise the \emph{boundary rule} used to validate proof
registries: a node may claim a threshold trust level only if every
non-dead-end child claims at least that level. The global rule is the
coKleisli form of the local check---the local predicate holds at every
position of the tree. We prove \emph{duplicate-stability}: if a tree
satisfies the global boundary rule, then every subtree does. Labels do
not disturb the navigation laws of the underlying directed container.
Two concrete trust chains (a seven-level agent chain and a six-level
loop-experiment chain) instantiate the general result. All stated
results have been machine-checked.
\end{abstract}

\tableofcontents
\bigskip

%======================================================================
\section{Introduction}
%======================================================================

Proof registries may be organised as finitely branching trees whose
nodes carry trust labels. A validator walks the tree and enforces a
node-local rule: a claim graded at or above a threshold may stand only
if every live child is graded at or above that threshold. Dead ends
(abandoned approaches) are exempt. Shared nodes import subtrees across
registry links; for validation to be compositional, a validated tree
must contribute only validated subtrees.

Directed containers supply the navigation structure: shapes are trees,
positions are paths from the root, and the comonad operations
$\mathrm{extract}$ and $\mathrm{duplicate}$ arise from path
concatenation. The boundary rule, stated at every position produced by
$\mathrm{duplicate}$, is stable under taking subtrees---transport along
the directed-container law that subtrees nest along path concatenation.

\paragraph{What this note claims.}
In ordinary mathematical language: validating the whole tree validates
every subtree. More precisely, if the local boundary check holds at
every position of a labelled tree~$t$, then it holds at every position
of every subtree $t \mathbin{\downarrow} p$. The proof is a one-line
transport along the nesting law for subtrees. Dead ends are modelled by
a boolean flag outside the trust order, because an order relation cannot
express an exemption. Two deployed trust chains instantiate the abstract
threshold.

%======================================================================
\section{Trust labels and labelled trees}
%======================================================================

\begin{definition}[Trust label]\label{def:label}
Let $T$ be a type of trust levels. A \textbf{trust label} is a pair
\begin{equation}
  \ell \;=\; \langle \mathrm{trust}(\ell),\; \mathrm{dead}(\ell) \rangle
  \;\in\; T \times \mathbb{B},
\end{equation}
consisting of a level in~$T$ and a dead-end flag. Dead ends lie outside
the trust order by design: the boundary rule \emph{exempts} them, and
no order relation expresses an exemption.
\end{definition}

\begin{definition}[Labelled rose tree]\label{def:ltree}
A \textbf{trust-labelled rose tree} over~$T$ is a finitely branching tree
with a label at every node:
\begin{equation}
  t \;=\; \mathrm{node}\bigl(\ell,\; n,\; (t_i)_{i < n}\bigr),
\end{equation}
where $\ell$ is a trust label, $n \in \mathbb{N}$ is the branching
factor, and each $t_i$ is a labelled rose tree. A leaf is a node with
$n = 0$.
\end{definition}

\begin{definition}[Root label]\label{def:lab}
The \textbf{root label} of $t = \mathrm{node}(\ell, n, (t_i))$ is
$\mathrm{lab}(t) = \ell$. Under the comonad structure below,
$\mathrm{lab}(t \mathbin{\downarrow} p)$ reads the label at position~$p$.
\end{definition}

%======================================================================
\section{Positions, subtrees, and path concatenation}
%======================================================================

\begin{definition}[Position]\label{def:lpos}
A \textbf{position} (path) in a labelled tree~$t$ is generated by
\begin{itemize}
  \item $\mathsf{here}$ --- the root of~$t$;
  \item $\mathsf{child}(i, p)$ --- descend into the $i$-th child and
        continue along path~$p$.
\end{itemize}
For a leaf ($n = 0$), the only position is $\mathsf{here}$.
\end{definition}

\begin{definition}[Subtree and shift]\label{def:sub-shift}
For a tree~$t$ and position~$p$ in~$t$:
\begin{itemize}
  \item the \textbf{subtree} $t \mathbin{\downarrow} p$ is~$t$ itself at
        $\mathsf{here}$, and the recursively computed subtree of the
        $i$-th child at $\mathsf{child}(i, p)$;
  \item the \textbf{shift} (path concatenation)
        $p \oplus q$ embeds a position~$q$ in the subtree at~$p$ back
        into~$t$: at $\mathsf{here}$ it is~$q$, and at
        $\mathsf{child}(i, p)$ it is $\mathsf{child}(i,\, p \oplus q)$.
\end{itemize}
\end{definition}

\begin{proposition}[Directed-container laws]\label{prop:d-laws}
Writing $o(t) = \mathsf{here}$ for the root position, the following hold
for all trees and positions (labels play no role):
\begin{enumerate}
  \item[\textbf{D1.}] $t \mathbin{\downarrow} o(t) \;=\; t$.
  \item[\textbf{D2.}] $o(t) \oplus q \;=\; q$.
  \item[\textbf{D3.}] $p \oplus o(t \mathbin{\downarrow} p) \;=\; p$.
  \item[\textbf{D4.}]
        $t \mathbin{\downarrow} (p \oplus q)
         \;=\;
         (t \mathbin{\downarrow} p) \mathbin{\downarrow} q$.
  \item[\textbf{D5.}]
        $(p \oplus q) \oplus r \;=\; p \oplus (q \oplus r)$
        (up to the identification of fibres supplied by~D4).
\end{enumerate}
In particular, labelled trees form a directed container; the navigation
structure---and hence the induced comonad---is independent of the labels.
\end{proposition}

\begin{proof}
By lean.
\end{proof}

\begin{remark}[Comonad reading]
The induced comonad on the extension has
\begin{align}
  \mathrm{extract}\,\langle t, v \rangle
  &\;=\; v(\mathsf{here}), \\
  \mathrm{duplicate}\,\langle t, v \rangle
  &\;=\;
  \bigl\langle
    t,\;
    p \mapsto
    \bigl\langle
      t \mathbin{\downarrow} p,\;
      q \mapsto v(p \oplus q)
    \bigr\rangle
  \bigr\rangle.
\end{align}
At every node, $\mathrm{duplicate}$ replaces the label with the full
sub-search from that point.
\end{remark}

%======================================================================
\section{The boundary rule}
%======================================================================

We now assume $T$ carries a partial order~$\le$ (in applications, a
linear chain of check strengths with bottom element).

\begin{definition}[Local boundary rule]\label{def:localok}
Fix a threshold $\theta \in T$. The \textbf{local boundary rule} at a
tree $t = \mathrm{node}(\ell, n, (t_i))$ is the proposition
\begin{equation}
  \localOk_\theta(t)
  \;:\!\!\Longleftrightarrow\;
  \Bigl(
    \theta \le \mathrm{trust}(\ell)
    \;\Longrightarrow\;
    \forall\, i < n.\;
    \bigl(
      \mathrm{dead}(\mathrm{lab}(t_i))
      \;\lor\;
      \theta \le \mathrm{trust}(\mathrm{lab}(t_i))
    \bigr)
  \Bigr).
\end{equation}
If the node claims $\theta$ or above, every non-dead child claims $\theta$
or above. There is no recursion: the check inspects one node and its
immediate children.
\end{definition}

\begin{definition}[Global boundary rule]\label{def:ok}
The \textbf{global boundary rule} at threshold~$\theta$ is the
coKleisli form of the local check: the local rule holds at
\emph{every} position,
\begin{equation}
  \Ok_\theta(t)
  \;:\!\!\Longleftrightarrow\;
  \forall\, p \text{ a position in } t.\;
  \localOk_\theta\bigl(t \mathbin{\downarrow} p\bigr).
\end{equation}
Equivalently: the local predicate holds on every sub-search produced by
$\mathrm{duplicate}$.
\end{definition}

%======================================================================
\section{Duplicate-stability}
%======================================================================

\begin{theorem}[Duplicate-stability of the trust boundary]
\label{thm:ok-sub}
Let $\theta \in T$ and let $t$ be a labelled tree. If $\Ok_\theta(t)$,
then for every position~$p$ in~$t$,
\begin{equation}
  \Ok_\theta\bigl(t \mathbin{\downarrow} p\bigr).
\end{equation}
\end{theorem}

\begin{proof}
By lean.
\end{proof}

\begin{remark}[Idea of the argument]
Positions of the subtree $t \mathbin{\downarrow} p$ embed into positions
of~$t$ by path concatenation $q \mapsto p \oplus q$. The nesting law~D4
identifies
\[
  (t \mathbin{\downarrow} p) \mathbin{\downarrow} q
  \;=\;
  t \mathbin{\downarrow} (p \oplus q),
\]
so the global check at $p \oplus q$ \emph{is} the check at~$q$ in the
subtree. Transport along this equality yields the result.
\end{remark}

\begin{theorem}[Whole equals all duplicates]\label{thm:ok-iff}
For every threshold~$\theta$ and tree~$t$,
\begin{equation}
  \Ok_\theta(t)
  \;\Longleftrightarrow\;
  \forall\, p \text{ a position in } t.\;
  \Ok_\theta\bigl(t \mathbin{\downarrow} p\bigr).
\end{equation}
The forward direction is Theorem~\ref{thm:ok-sub}; the backward direction
reads the root component of each duplicate.
\end{theorem}

\begin{proof}
By lean.
\end{proof}

\begin{remark}
This is the load-bearing fact for shared nodes: a stub pointing into a
validated registry inherits a validated subtree, and the validator never
needs to re-check across the link.
\end{remark}

%======================================================================
\section{Deployed trust chains}
%======================================================================

\begin{definition}[Agent trust chain]\label{def:trustlevel}
The \textbf{agent trust chain} is the seven-element linear order
\begin{align*}
  &\mathrm{unchecked}
  \;<\; \mathrm{speculative}
  \;<\; \mathrm{computed}
  \;<\; \mathrm{proved} \\
  &\qquad
  \;<\; \mathrm{peerReviewed}
  \;<\; \mathrm{published}
  \;<\; \mathrm{leanVerified},
\end{align*}
with bottom element $\mathrm{unchecked}$. The dead-end status is not a
chain element: it is the $\mathrm{dead}$ flag on labels.
\end{definition}

\begin{definition}[Registry validity]\label{def:registryok}
A labelled tree over the agent trust chain is \textbf{registry-valid} if
it satisfies the global boundary rule at threshold $\mathrm{proved}$:
\begin{equation}
  \mathrm{RegistryOk}(t)
  \;:\!\!\Longleftrightarrow\;
  \Ok_{\mathrm{proved}}(t).
\end{equation}
\end{definition}

\begin{corollary}[Deployed duplicate-stability]\label{cor:registryok-sub}
If $t$ is registry-valid, then every subtree of~$t$ is registry-valid.
\end{corollary}

\begin{proof}
By lean.
\end{proof}

\begin{definition}[Loop-experiment trust chain]\label{def:looptrust}
The \textbf{loop-experiment trust chain} is the six-element linear order
\begin{align*}
  &\mathrm{unchecked}
  \;<\; \mathrm{speculative}
  \;<\; \mathrm{computed}
  \;<\; \mathrm{proved} \\
  &\qquad
  \;<\; \mathrm{verified}
  \;<\; \mathrm{leanVerified},
\end{align*}
with bottom element $\mathrm{unchecked}$. The grade $\mathrm{verified}$
replaces the social-process grades $\mathrm{peerReviewed}$ and
$\mathrm{published}$ of Definition~\ref{def:trustlevel}: it records an
independent hostile review.
\end{definition}

\begin{corollary}[Loop duplicate-stability]\label{cor:loopok-sub}
A labelled tree over the loop-experiment chain that satisfies
$\Ok_{\mathrm{proved}}$ has every subtree satisfying
$\Ok_{\mathrm{proved}}$.
\end{corollary}

\begin{proof}
By lean.
\end{proof}

\begin{remark}
Both corollaries are instances of Theorem~\ref{thm:ok-sub}, which is
parametric in the trust type~$T$ and the threshold~$\theta$. The same
structural fact applies to any chain and any threshold.
\end{remark}

%======================================================================
\section{Discussion}
%======================================================================

\paragraph{What has been established.}
Trust-labelled rose trees form a directed container whose navigation is
independent of labels. The boundary rule---local at each node, global
as a coKleisli condition---is stable under $\mathrm{duplicate}$:
validated trees decompose into validated subtrees. Two concrete chains
used in deployed registries instantiate the general theorem.

\paragraph{Relation to transport of trust.}
Duplicate-stability says validated trees decompose. Separately, morphisms
of trust systems transport canonical trust laxly: importing along a
morphism never claims more than the image of the source's canonical
trust. Together these are the two halves of compositional registry
validation: structure decomposes, and grades move only downward under
import.

\paragraph{Machine checking.}
Every definition and theorem statement above corresponds to a declaration
in a machine-checked formalisation. The phrase ``by lean'' indicates that
the proof has been verified by a proof assistant.

%======================================================================
\section*{Acknowledgements}
%======================================================================

The presentation here is a prose rendering of a formal development
intended for human review of the theorem statements. The directed
container laws follow the theory of Ahman, Chapman, and Uustalu; the
boundary rule is the mathematical content of a node-local registry
validator.

\begin{thebibliography}{9}

\bibitem{ahman2012}
D.~Ahman, J.~Chapman, and T.~Uustalu.
\newblock When is a container a comonad?
\newblock In \emph{FoSSaCS}, volume~7213 of \emph{LNCS}, pages 74--88.
\newblock Springer, 2012.

\bibitem{ahman2014}
D.~Ahman and T.~Uustalu.
\newblock Directed containers as categories.
\newblock \emph{EPTCS}, 207:89--98, 2016.

\end{thebibliography}

\end{document}
